\chapter{Images as Two-Dimensional Signals}

For thirty chapters we have lived on a line. A signal was a function of one
variable -- almost always time -- and everything we built, from sinusoids to the
Fourier transform to convolution, was a story about how a single number wiggles
as that one variable marches forward. Now we are going to add an axis. An image
is what you get when the thing you measure depends not on \emph{when} but on
\emph{where}: not $x(t)$ but $f(x,y)$, a brightness that varies as you move left
and right \emph{and} up and down across a plane.

This is a bigger leap than it sounds, and a smaller one than you fear. Bigger,
because two dimensions buy you genuinely new phenomena -- edges have
orientation, textures have direction, and ``frequency'' splits into a vector that
points somewhere. Smaller, because almost every tool you already own transfers
across with one extra index bolted on. A pixel grid is just a 2-D array of
samples. Brightness is just amplitude. The histogram you may have met for a
random variable is exactly the right object for understanding contrast. And in
the next chapters the 2-D Fourier transform will turn out to be the 1-D
transform applied twice, once along each axis.

This chapter lays the foundation. We will see an image as a sampled function on a
grid, understand what a pixel actually is and what intensity and color mean,
think carefully about 2-D sampling and resolution (the same Nyquist story, now in
a plane), meet brightness, contrast, and the all-important histogram, and learn
the simplest family of image operations -- \emph{point operations}, which touch
each pixel in isolation. By the end you will look at any image and see a signal.

\chapterline

\section{From one axis to two}
\label{ch30:one-to-two}

Start with what you know. A discrete-time signal is a list of numbers indexed by
an integer: $x[n]$, where $n$ runs $0,1,2,\dots$ and each $x[n]$ is the value of
the signal at that sample. Plot it and you get the familiar stem-and-line
picture: time across, amplitude up.

A grayscale image is the same idea with \emph{two} integer indices. We write
$f[m,n]$, where $m$ selects a row and $n$ selects a column, and $f[m,n]$ is the
brightness of the picture at that grid location. Instead of a list you have a
\emph{rectangular array} of numbers -- a matrix, in the linear-algebra sense.
Where a 1-D signal of length $N$ is a point in $\R^N$, an image of height $M$ and
width $N$ is a point in $\R^{M\times N}$: a single vector with $MN$ coordinates,
one per pixel.

\begin{intuitionbox}
Picture the difference physically. A 1-D signal is a single microphone reporting
one pressure value as time passes. A 2-D image is a \emph{room full} of tiny
light meters arranged in a grid, each reporting how bright it is at its own spot,
all at the same instant. Time has been replaced by position. The signal no longer
flows past you; it sits there, spread across a plane, and you read it by moving
your eye instead of waiting.
\end{intuitionbox}

It helps to keep both the \emph{continuous} and the \emph{sampled} pictures in
mind, exactly as we did for time signals. The world presents a continuous image
$f(x,y)$ -- the real light field falling on a sensor, defined for every real
position $(x,y)$ in some rectangle. A camera cannot store a function of real
numbers, so it \emph{samples}: it measures the light only on a grid of points and
\emph{quantizes} each measurement to a finite number of levels. The result is the
array $f[m,n]$. The continuous-to-discrete jump is the same one from the sampling
chapter, just performed along two spatial axes instead of one temporal axis.

\begin{figure}[H]
\centering
\begin{tikzpicture}[scale=1.0]
% continuous "image" as a smooth shaded gradient
\begin{scope}
\foreach \i in {0,...,19}{
  \pgfmathsetmacro{\g}{40+45*sin(deg(0.45*\i))*sin(deg(0.45*\i))}
  \definecolor{cc}{gray}{\g/100}
  \fill[cc] (0.2*\i,0) rectangle (0.2*\i+0.2,2);
}
\draw[thick] (0,0) rectangle (4,2);
\node[below] at (2,-0.15) {continuous $f(x,y)$};
\end{scope}
% arrow
\draw[-{Latex[length=3mm]}, thick] (4.4,1) -- (5.4,1) node[midway,above,font=\small]{sample};
% sampled grid of squares
\begin{scope}[shift={(5.8,0)}]
\foreach \c in {0,...,7}{
  \pgfmathsetmacro{\g}{40+45*sin(deg(0.45*2.5*\c))*sin(deg(0.45*2.5*\c))}
  \definecolor{cc}{gray}{\g/100}
  \foreach \r in {0,...,3}{
    \fill[cc] (0.5*\c,0.5*\r) rectangle (0.5*\c+0.5,0.5*\r+0.5);
    \draw[gray!50] (0.5*\c,0.5*\r) rectangle (0.5*\c+0.5,0.5*\r+0.5);
  }
}
\draw[thick] (0,0) rectangle (4,2);
\node[below] at (2,-0.15) {sampled $f[m,n]$};
\end{scope}
\end{tikzpicture}
\caption{The continuous-to-discrete jump in 2-D. The real light field (left) is a
function on a continuous rectangle; the camera samples it on a grid and stores
one number per cell (right). Each cell is a pixel.}
\label{ch30:fig-sample}
\end{figure}

\begin{focusbox}
Two conventions will save you grief. \textbf{(1)~Row, then column.} We index
$f[m,n]$ with $m$ as the row (vertical) and $n$ as the column (horizontal), like
a matrix. \textbf{(2)~Rows count downward.} In image coordinates $m=0$ is the
\emph{top} row and increases as you go down, because that is the order pixels
stream out of a sensor and into a file. This flips the $y$-axis relative to a
math plot, and is the source of countless upside-down images. We will be explicit
whenever it matters.
\end{focusbox}

\section{What a pixel really is}
\label{ch30:pixel}

The word ``pixel'' (picture element) gets used for two different things, and
conflating them causes real bugs. A pixel is, first, a \emph{sample}: a single
number (or small group of numbers) measured at one grid location. A pixel is
also, loosely, the little \emph{square} you see when you zoom in on a screen.
Those squares are a \emph{rendering choice}, not the data. The data is a set of
point samples; the squares are how we paint each sample so the eye can see it.

\begin{trapbox}
``A pixel is a tiny square.'' This is the most persistent misconception in
imaging. The samples are points; the square (or, on a sensor, the little
light-collecting area) is the reconstruction. Believing pixels are squares leads
people to expect that ``zooming in'' should show sharper squares, when in truth a
correct reconstruction should show a smoothly interpolated surface -- the 2-D
analog of connecting the dots in a 1-D plot. Alvy Ray Smith's classic 1995 note
``A Pixel Is Not A Little Square'' is the canonical takedown.
\end{trapbox}

A grayscale pixel holds one number: its \textbf{intensity} or brightness. How
that number is stored is the matter of \emph{quantization}, which we treat
properly in the chapter on quantization; here we only need the headline. Most
photographs store intensity in 8 bits, giving integer levels $0,1,\dots,255$,
where $0$ is black and $255$ is white. Eight bits is not sacred -- medical and
scientific images often use 12, 14, or 16 bits, and internal processing is
frequently done in floating point on the range $[0,1]$ to avoid rounding. The
key idea is that the stored value is a quantized estimate of a continuous
physical brightness.

\begin{connectionbox}
The image-as-vector view is the bridge to the linear-algebra companion. An
$M\times N$ image \emph{is} a vector in $\R^{MN}$; flattening it row by row (or
column by column) is just choosing coordinates. Every operation we will study --
brightening, blurring, edge detection -- is a function on this vector space, and
the nicest ones are \emph{linear}, i.e.\ matrices. The inner product
$\inner{f}{g}=\sum_{m,n} f[m,n]\,g[m,n]$ measures how much two images overlap,
exactly the $\inner{\cdot}{\cdot}$ we used to project signals onto sinusoids in
the signals-as-vectors chapter. Orthogonality, projection, and bases all carry
over unchanged.
\end{connectionbox}

\subsection{Color: three signals stacked}
\label{ch30:color}

A color image is not one signal but several, sampled on the same grid. The human
eye has three kinds of color-sensitive cone, roughly tuned to long, medium, and
short wavelengths, so three numbers per location are enough to reproduce almost
any color a person can see. The default representation is \textbf{RGB}: three
arrays $R[m,n]$, $G[m,n]$, $B[m,n]$ giving the red, green, and blue intensities.
Stacked, they form a \emph{three-channel} image -- think of it as three grayscale
signals lying on top of one another, or as a single function whose value is a
vector $(R,G,B)$ rather than a scalar.

\begin{happeningbox}
What is really happening with color is that we have moved from a scalar-valued
signal to a \emph{vector-valued} one. Each pixel carries a little vector in a
3-D ``color space''. Nothing about the geometry of the grid changes -- the
sampling, the resolution, the neighbor relationships are all identical to the
grayscale case. We have simply attached a 3-vector to every grid point. Every
single-channel technique in this book applies to a color image by running it on
each channel (with the caveat that the channels are correlated, so independent
processing can shift colors).
\end{happeningbox}

RGB is convenient for displays (which literally have red, green, and blue
emitters) but awkward for analysis, because brightness and color are tangled
across all three numbers. Make something twice as bright and all three of $R$,
$G$, $B$ change. So we often convert to a space that \emph{separates} a single
brightness channel from the color. The most useful for our purposes is some form
of \textbf{luma plus chroma}. A standard luma (the brightness a grayscale viewer
would see) is the weighted sum
\begin{equation}
\label{ch30:luma}
Y = 0.299\,R + 0.587\,G + 0.114\,B,
\end{equation}
with green weighted most because the eye is most sensitive to green. The unequal
weights are not arbitrary: they are the eye's own sensitivity curve, baked into
the standard. The remaining two numbers (the chroma) describe hue and
saturation. Closely related spaces include \textbf{HSV}/\textbf{HSL} (hue,
saturation, value/lightness), which are handy when you want to threshold ``how
red is this'' independently of ``how bright is this.''

\begin{table}[H]
\centering
\caption{Common color representations and what each is good for.}
\label{ch30:tab-color}
\begin{tabular}{@{}lll@{}}
\toprule
Space & Components & Good for \\
\midrule
RGB & red, green, blue & displays, raw capture, storage \\
Grayscale & single luma $Y$ & analysis where color is irrelevant \\
YCbCr & luma + 2 chroma & compression (JPEG), video \\
HSV/HSL & hue, saturation, value & color thresholding, painting tools \\
\bottomrule
\end{tabular}
\end{table}

\begin{whybox}
Why bother converting to luma before analysis? Because most of the
\emph{structure} we care about -- edges, shapes, motion -- lives in brightness,
not color. The eye agrees: it has far more brightness-sensitive rods than
color-sensitive cones, and far sharper spatial acuity for luma than for chroma.
This is exactly why JPEG and every video codec store luma at full resolution but
\emph{subsample} the chroma (the ``4:2:0'' you see in codec settings): you can
halve the color resolution and almost nobody notices. Working in luma both
matches the data's information content and cuts your computation by a factor of
three.
\end{whybox}

\section{Sampling and resolution in two dimensions}
\label{ch30:sampling}

The number of samples is the image's \textbf{resolution}: an image that is
$1920\times 1080$ has $1920$ columns and $1080$ rows, about two million pixels.
But raw pixel count is only half the story. What matters for fidelity is the
\emph{sample spacing} relative to how fast the scene's brightness changes -- the
exact 2-D version of Nyquist.

Recall the 1-D rule: to capture a signal without aliasing, you must sample at
more than twice the highest frequency present. In 2-D the same rule applies
\emph{separately along each axis}, in terms of \textbf{spatial frequency} -- how
many brightness cycles fit per unit distance. A picket fence photographed
straight on is a near-sinusoid in the horizontal direction; if its bars are
finer than two samples per cycle, you get aliasing. Sample spacing $\Delta x$
lets you represent horizontal spatial frequencies up to
$1/(2\Delta x)$ cycles per unit length without aliasing; the same holds for
$\Delta y$ vertically.

\begin{figure}[H]
\centering
\begin{tikzpicture}
\begin{axis}[
  width=12cm, height=4.6cm,
  xlabel={horizontal position (pixels)}, ylabel={brightness},
  xmin=0, xmax=16, ymin=-0.2, ymax=1.2,
  axis lines=box, samples=200, domain=0:16,
  legend pos=outer north east, legend style={font=\footnotesize},
  title={A fine pattern that the grid cannot tell from a coarse one},
]
\addplot[specred, thick, domain=0:16, samples=200] {0.5+0.45*cos(deg(2*pi*0.75*x))};
\addlegendentry{true pattern ($0.75$ cyc/px)}
\addplot[sigblue, dashed, thick, domain=0:16, samples=200] {0.5+0.45*cos(deg(2*pi*0.25*x))};
\addlegendentry{alias ($0.25$ cyc/px)}
\addplot[black, only marks, mark=*, mark size=1.4pt, domain=0:16, samples=17]
  {0.5+0.45*cos(deg(2*pi*0.75*x))};
\addlegendentry{samples (1 per px)}
\end{axis}
\end{tikzpicture}
\caption{One row of pixels through a fine periodic pattern. The true pattern (red,
$0.75$ cycles per pixel) is finer than the Nyquist limit of $0.5$ cycles per pixel,
so the samples (dots, one per pixel) land in exactly the same places as a much
coarser pattern (blue dashed, $0.25$ cycles per pixel). The grid cannot distinguish
the two: the fine pattern masquerades as the coarse one. This is aliasing, exactly
as in 1-D, now along a spatial axis.}
\label{ch30:fig-alias1d}
\end{figure}

\begin{happeningbox}
The famous \emph{moir\'e} patterns you see when a camera photographs a striped
shirt or a brick wall, and the jagged ``stair-steps'' on diagonal lines, are both
spatial aliasing. High spatial frequencies in the scene -- finer than two samples
per cycle -- fold down and appear as spurious low-frequency patterns that were
never in the scene. Camera designers fight this with an \emph{optical low-pass
filter} (a slight deliberate blur in front of the sensor) -- the optical version
of the anti-aliasing filter you put before an ADC, and the first stage of the
capture pipeline in Figure~\ref{ch30:fig-pipeline}.
\end{happeningbox}

Resolution is therefore a contract between sample spacing and the finest detail
you intend to capture. More pixels help only up to the point where the optics
(and the anti-alias filter) stop delivering finer detail; past that you are
sampling blur. This is why a 50-megapixel phone photo through a smudged plastic
lens can carry less real information than a sharp 12-megapixel shot.

\begin{figure}[H]
\centering
\begin{tikzpicture}[
  node distance=6mm and 9mm,
  block/.style={draw, rounded corners, align=center, font=\small,
                minimum height=10mm, minimum width=17mm, fill=blue!4},
  >={Latex[length=2.5mm]}]
\node[block] (scene) {scene\\$f(x,y)$};
\node[block, right=of scene] (lens) {lens +\\optical\\low-pass};
\node[block, right=of lens] (samp) {sample\\on grid};
\node[block, right=of samp] (quant) {quantize\\(8--16 bit)};
\node[block, right=of quant] (arr) {array\\$f[m,n]$};
\draw[->] (scene) -- (lens);
\draw[->] (lens) -- node[above,font=\footnotesize]{anti-alias} (samp);
\draw[->] (samp) -- (quant);
\draw[->] (quant) -- (arr);
\end{tikzpicture}
\caption{The capture pipeline, stage for stage the 2-D twin of the 1-D
analog-to-digital chain. The optical low-pass filter band-limits the light field
\emph{before} the grid samples it -- the spatial anti-alias filter -- and
quantization rounds each sample to a finite set of levels.}
\label{ch30:fig-pipeline}
\end{figure}

\begin{connectionbox}
The same sampling theorem governs the time axis for our IMU and EMG signals
(the sampling-and-Nyquist chapter) and the spatial axes here. ``Sample fast
enough relative to the fastest variation'' is one theorem wearing two costumes.
The difference is dimensionality: 1-D signals alias into a wrong tone, 2-D
signals alias into a wrong \emph{pattern with an orientation}.
\end{connectionbox}

\section{Brightness and contrast}
\label{ch30:brightness-contrast}

Two of the most basic image qualities have precise meanings as operations on the
pixel values. \textbf{Brightness} is the overall level -- the average intensity.
Add a constant to every pixel and you raise the brightness, shifting the whole
image lighter (or darker, for a negative constant):
\begin{equation}
\label{ch30:brightness}
g[m,n] = f[m,n] + b .
\end{equation}
\textbf{Contrast} is the \emph{spread} of intensities -- how far the bright and
dark pixels sit from the middle. Multiply every pixel's deviation from a midpoint
by a gain $a>1$ and you stretch the values apart, increasing contrast; $a<1$
squeezes them together, washing the image out:
\begin{equation}
\label{ch30:contrast}
g[m,n] = a\,\bigl(f[m,n]-c\bigr) + c,
\end{equation}
where $c$ is the pivot (often the mid-gray $127.5$ for 8-bit data). Combine the
two and you have the general affine intensity map $g = a\,f + b'$.

\begin{intuitionbox}
Think of the intensities as dots on a number line from black to white.
\emph{Brightness} slides all the dots together left or right. \emph{Contrast}
pulls them apart from (or pushes them toward) a center, like stretching a rubber
band pinned at mid-gray. Neither operation cares \emph{where} a pixel is in the
image -- only what value it holds. That locality-blindness is the defining
feature of the point operations we study next.
\end{intuitionbox}

\begin{trapbox}
After an affine map some results fall outside the valid range -- below $0$ or
above $255$. You must \emph{clip} (clamp) them back into range. Forgetting to clip,
or letting an 8-bit integer wrap around (so $256$ becomes $0$), produces ugly
artifacts: blown highlights turning black, or speckle where bright areas
overflow. Always saturate, never wrap.
\end{trapbox}

\section{The histogram: a signal's value distribution}
\label{ch30:histogram}

If you want to understand an image's brightness and contrast at a glance, do not
stare at the pixels -- look at the \textbf{histogram}. The histogram counts how
many pixels hold each intensity value. For 8-bit data it is an array $h[k]$,
$k=0,\dots,255$, where $h[k]$ is the number of pixels with intensity exactly $k$.
Normalize by the total pixel count $MN$ and you get $p[k]=h[k]/(MN)$, the
\emph{fraction} of pixels at each level -- a probability distribution over
intensity.

\begin{figure}[H]
\centering
\begin{tikzpicture}
\begin{axis}[
  width=12cm, height=4.8cm,
  xlabel={intensity $k$ (0 = black, 255 = white)},
  ylabel={fraction of pixels $p[k]$},
  xmin=0, xmax=255, ymin=0, ymax=0.018,
  axis lines=box, ytick=\empty,
  legend pos=north east, legend style={font=\footnotesize},
]
\addplot[sigblue, thick, samples=120, domain=0:255]
  {0.016*exp(-((x-128)^2)/(2*22^2))};
\addlegendentry{low contrast}
\addplot[targetgreen, thick, samples=120, domain=0:255]
  {0.006*exp(-((x-90)^2)/(2*60^2))+0.005*exp(-((x-180)^2)/(2*45^2))};
\addlegendentry{full range}
\end{axis}
\end{tikzpicture}
\caption{Two histograms. The narrow blue distribution uses only a sliver of the
available range -- a flat, low-contrast image. The wide green distribution spreads
across the range and is bimodal (a dark subject on a light background, say) -- a
punchy, well-exposed image.}
\label{ch30:fig-hist}
\end{figure}

The shape of the histogram tells you nearly everything about exposure and
contrast. A histogram bunched near $0$ means an underexposed (dark) image; bunched
near $255$ means overexposed. A narrow spike anywhere means low contrast -- the
pixels barely differ. A distribution spread across the whole axis means good
contrast. Two separated humps often mean a clear subject and background.
Crucially, the histogram \emph{throws away all spatial information}: shuffle the
pixels into random positions and the histogram is unchanged. It is purely a
statement about the distribution of \emph{values}.

\begin{connectionbox}
The normalized histogram $p[k]$ is a genuine probability mass function over the
random variable ``intensity of a pixel picked at random,'' which is exactly the
object from the probability companion. Its mean
$\bar{f}=\sum_k k\,p[k]=\E[\text{intensity}]$ is the image's average brightness,
and its variance $\sum_k (k-\bar f)^2\,p[k]=\Var[\text{intensity}]$ is a clean
measure of contrast. Brightness and contrast, defined informally above, are just
the first two moments of this distribution. Histogram equalization (below) is the
``make it uniform'' transform from probability, applied to pixels.
\end{connectionbox}

\subsection{Histogram stretching and equalization}
\label{ch30:equalization}

If the histogram is narrow, you can improve contrast by \textbf{stretching} it to
fill the range. Find the actual minimum and maximum intensities $f_{\min}$ and
$f_{\max}$ in the image and apply the linear map that sends $f_{\min}\to 0$ and
$f_{\max}\to 255$:
\begin{equation}
\label{ch30:stretch}
g[m,n] = 255\cdot \frac{f[m,n]-f_{\min}}{f_{\max}-f_{\min}}.
\end{equation}
This is just the affine map of \eqref{ch30:contrast} with the gain and pivot
chosen automatically from the data.

\textbf{Histogram equalization} is cleverer. Instead of a linear stretch, it
applies a nonlinear remapping that makes the \emph{output} histogram as flat
(uniform) as possible, allocating the most output range to the most populated
input levels. The recipe is the cumulative distribution: compute the running
total
\begin{equation}
\label{ch30:cdf}
C[k] = \sum_{i=0}^{k} p[i],
\end{equation}
then map each intensity $k\mapsto \floor{255\,C[k]}$. Pixel values that are common
get spread far apart (high contrast where it matters); rare values get squeezed
together.

\begin{figure}[H]
\centering
\begin{tikzpicture}
\begin{axis}[
  width=10.5cm, height=5.6cm,
  xlabel={input intensity $k$}, ylabel={output $\,\floor{255\,C[k]}$},
  xmin=0, xmax=160, ymin=0, ymax=270,
  axis lines=box, legend pos=south east, legend style={font=\footnotesize},
  title={The CDF \emph{is} the equalization transfer curve},
]
% scaled CDF staircase for the worked-example image
\addplot[sourceblue, thick, const plot mark left] coordinates {
  (30,28) (60,113) (90,198) (120,226) (150,255)};
\addlegendentry{$255\,C[k]$ (equalization map)}
\addplot[sourceblue, only marks, mark=*, mark size=1.4pt] coordinates {
  (30,28) (60,113) (90,198) (120,226) (150,255)};
\addplot[gray, dashed] coordinates {(0,0) (160,255)};
\addlegendentry{linear stretch (for comparison)}
\end{axis}
\end{tikzpicture}
\caption{Histogram equalization for the worked-example patch of
Section~\ref{ch30:worked}. The transfer curve is the scaled CDF: it climbs
steeply where many pixels live (the crowded $60$ and $90$ levels) and gently where
pixels are sparse. The steep segments stretch common intensities far apart, which
is where the contrast goes. The dashed line is the linear stretch, which ignores
how many pixels sit at each level.}
\label{ch30:fig-cdf}
\end{figure}

\begin{whybox}
Why does mapping through the CDF flatten the histogram? Because the CDF is exactly
the transform that turns \emph{any} distribution into a uniform one -- the
``probability integral transform'' from the probability book. If a random
variable $X$ has CDF $C$, then $C(X)$ is uniform on $[0,1]$. Applying $C$ to the
pixel values therefore drives the value distribution toward flat, which means
intensities are now spread evenly across the available range -- the definition of
maximum global contrast. (In practice it is approximate, because intensities are
discrete and many pixels share a value, so the result is uniform-ish, not exactly
uniform.)
\end{whybox}

\begin{labbox}
With any image library, load a photo, convert to grayscale, and plot its
histogram. Then try the three contrast operations -- a manual affine map, a
min/max stretch, and equalization -- and re-plot the histogram each time. Watch
how stretching widens the support and equalization flattens it. Do it on a
deliberately underexposed photo to see the dramatic version.
\end{labbox}

\begin{lstlisting}[language=Python,caption={Grayscale conversion, histogram, and equalization in NumPy.},captionpos=b]
import numpy as np

def to_gray(rgb):                          # rgb: H x W x 3, uint8
    w = np.array([0.299, 0.587, 0.114])
    return rgb.astype(float) @ w           # H x W, float luma

def histogram(gray, levels=256):
    h, _ = np.histogram(gray, bins=levels, range=(0, levels))
    return h

def equalize(gray):
    g = np.clip(gray, 0, 255).astype(int)
    h = histogram(g)
    cdf = np.cumsum(h) / g.size            # normalized CDF in [0,1]
    lut = np.floor(255 * cdf).astype(np.uint8)   # 256-entry lookup table
    return lut[g]                          # remap every pixel via the LUT
\end{lstlisting}

\section{Point operations: one pixel at a time}
\label{ch30:point-ops}

Brightness, contrast, stretching, and equalization are all examples of a
\textbf{point operation} (also called a \emph{pixel} or \emph{intensity}
transform): a rule that computes each output pixel from \emph{only} the input
pixel at the same location,
\begin{equation}
\label{ch30:point-op}
g[m,n] = T\bigl(f[m,n]\bigr),
\end{equation}
where $T$ is a single-input function $T:\R\to\R$ -- the same $T$ for every
location. Because $T$ ignores the neighbors entirely, a point operation can be
implemented as a $256$-entry \textbf{lookup table} (LUT) for 8-bit images:
precompute $T(k)$ for $k=0,\dots,255$ once, then replace each pixel by a table
read. This is blisteringly fast and is how real-time brightness/contrast sliders
work.

\begin{figure}[H]
\centering
\begin{tikzpicture}
\begin{axis}[
  width=10.5cm, height=6.2cm,
  xlabel={input intensity $f$}, ylabel={output intensity $g=T(f)$},
  xmin=0, xmax=255, ymin=0, ymax=255,
  axis lines=box, samples=120, domain=0:255,
  legend pos=outer north east, legend style={font=\footnotesize},
  title={Transfer curves of common point operations},
]
\addplot[gray, dashed, domain=0:255, samples=2] {x};
\addlegendentry{identity}
\addplot[specred, thick, domain=0:255, samples=2] {255-x};
\addlegendentry{negative}
\addplot[sigblue, thick, domain=0:255, samples=120] {255*(x/255)^0.45};
\addlegendentry{gamma 0.45 (brighten)}
\addplot[targetgreen, thick, domain=0:255, samples=120] {255*(x/255)^2.2};
\addlegendentry{gamma 2.2 (darken)}
\addplot[quotientorange, thick, domain=0:255, samples=120]
  {255/(1+exp(-0.05*(x-128)))};
\addlegendentry{sigmoid (S-curve)}
\end{axis}
\end{tikzpicture}
\caption{A point operation is fully described by its \emph{transfer curve}
$g=T(f)$: input intensity on the horizontal axis, output on the vertical. The
diagonal is the do-nothing identity. Reflecting it gives the negative; bending it
up brightens; bending it down darkens; an S-shape boosts mid-tone contrast.}
\label{ch30:fig-transfer}
\end{figure}

The transfer curve is the single most useful way to think about point operations.
Reading it is a skill: where the curve is \emph{steep}, small input differences
become large output differences -- contrast is increased for that band of
intensities. Where the curve is \emph{flat}, input differences are crushed --
contrast lost. A straight diagonal preserves everything. With that lens, every
point operation becomes obvious from its shape.

\begin{example}[The photographic negative]
\label{ch30:ex-negative}
The negative is the point operation $T(f)=255-f$ (for 8-bit data). It maps black
to white and white to black, reflecting every value about mid-gray. Its transfer
curve is the descending diagonal in Figure~\ref{ch30:fig-transfer}. Because the
slope has magnitude $1$ everywhere, contrast is exactly preserved -- only the
sense of light and dark is flipped. This is why a film negative and its print
carry the same detail.
\end{example}

\subsection{Gamma: the most important nonlinear point operation}
\label{ch30:gamma}

Not all point operations are linear, and the most consequential one is
\textbf{gamma correction}:
\begin{equation}
\label{ch30:gamma-eq}
g = 255\left(\frac{f}{255}\right)^{\gamma}.
\end{equation}
For $\gamma<1$ the curve bows upward, brightening shadows and compressing
highlights; for $\gamma>1$ it bows downward, darkening. Gamma exists for two deep
reasons. The first is that displays are physically nonlinear: the light a CRT or
LCD emits is roughly the input value raised to a power near $2.2$, so image data is
\emph{pre-distorted} by the inverse to come out right. The second reason is the one
that makes gamma stick around: human brightness perception is itself roughly
logarithmic (Weber-Fechner), so we perceive equal \emph{ratios} of light as equal
steps. Encoding
images in a gamma-compressed space therefore spends bits where the eye can see
differences, packing more perceptual resolution into the dark tones we are
sensitive to.

\begin{historybox}
Gamma is a fossil of cathode-ray-tube physics. The electron-gun brightness of a
CRT followed a power law $L\propto V^{\gamma}$ with $\gamma\approx 2.2$ -- not a
choice, just how the hardware behaved. Early television engineers, rather than fix
it at the display, pre-corrected the signal at the camera with the inverse power.
By a happy accident this pre-correction also matched the eye's roughly logarithmic
response, so when flat-panel displays arrived the entire imaging pipeline kept the
$2.2$ gamma on purpose. A hardware quirk became a perceptual standard, and the
sRGB color space enshrines a curve very close to it.
\end{historybox}

\begin{trapbox}
Do your linear math in linear light. Blurring, resizing, and averaging assume the
pixel values add up the way light adds up -- but stored values are usually
\emph{gamma-encoded}, not linear. Averaging two gamma-encoded pixels gives the
wrong brightness (famously, a 50/50 blend of pure red and pure green comes out a
muddy dark instead of a bright yellow). For correct results, undo gamma to linear,
do the arithmetic, then re-apply gamma. Point operations like brightness are
usually fine in encoded space; \emph{neighbor} operations are not.
\end{trapbox}

\subsection{Thresholding}
\label{ch30:threshold}

The most aggressive point operation is the \textbf{threshold}, which collapses
the whole intensity axis to two values:
\begin{equation}
\label{ch30:threshold-eq}
g[m,n] = \begin{cases} 255 & f[m,n] \ge \theta,\\ 0 & f[m,n] < \theta.\end{cases}
\end{equation}
Its transfer curve is a single step at $\theta$. Thresholding turns a grayscale
image into a black-and-white \emph{mask} -- the foundation of segmentation,
document scanning, and barcode reading. Choosing $\theta$ well is the whole game;
a classic automatic choice is \emph{Otsu's method}, which picks the threshold that
best separates the histogram into two clusters (minimizing within-class variance),
once again treating the histogram as a probability distribution.

\begin{applicationbox}
Point operations are the front line of medical and scientific imaging. An X-ray
or CT slice arrives as 12- or 16-bit data spanning a vast intensity range; no
display and no eye can show it all at once. Radiologists apply a \emph{window/level}
point operation -- an affine map that linearly stretches one narrow band of
intensities (say, soft tissue, or bone, or lung) across the full display range and
clips the rest. It is brightness-and-contrast \eqref{ch30:contrast} with clinically
chosen $a$ and $c$, and switching ``windows'' is how one scan reveals completely
different anatomy.
\end{applicationbox}

\section{The jump from 1-D to 2-D, made precise}
\label{ch30:jump}

We have been leaning on the slogan ``an image is just a 2-D signal.'' It is worth
spelling out exactly what carries over for free and what is genuinely new, because
the next chapters depend on knowing the difference.

\textbf{What transfers unchanged.} Everything that treats the signal pointwise or
as a vector: amplitude, sampling and Nyquist, quantization, the histogram,
point operations, the inner product, energy ($\sum f^2$), and linearity. A 2-D
LTI system (we will call it shift-invariant) still acts by convolution -- now a
\emph{2-D} convolution that slides a small kernel over the grid. The Fourier
transform still exists and still diagonalizes convolution; it just becomes a
transform over two frequency variables, $u$ (horizontal) and $v$ (vertical).

\textbf{What is genuinely new.} Direction. On a line there are only two ways to
move (left, right) and ``frequency'' is a single scalar. On a plane you can move
in any direction, so spatial frequency becomes a \emph{vector} $(u,v)$: it has a
magnitude (how fine the pattern) \emph{and} an orientation (which way the stripes
run). A pattern of vertical stripes and a pattern of horizontal stripes have the
same fineness but live at different points in 2-D frequency space. Edges, too,
acquire orientation. This is why the next chapter's 2-D spectra are pictures, not
curves -- and why image processing has a richness that 1-D signal processing
lacks.

\begin{figure}[H]
\centering
\begin{tikzpicture}
% three little patches showing orientation of 2D spatial frequency
\foreach \k/\ang/\lab in {0/0/{vertical stripes}, 1/90/{horizontal stripes}, 2/45/{diagonal stripes}}{
  \begin{scope}[shift={(\k*4.2,0)}]
    \foreach \i in {0,...,23}{
      \pgfmathsetmacro{\px}{0.125*\i}
      \foreach \j in {0,...,23}{
        \pgfmathsetmacro{\py}{0.125*\j}
        \pgfmathsetmacro{\g}{50+45*cos(deg(2*pi*1.6*(\px*cos(\ang)+\py*sin(\ang))))}
        \definecolor{cc}{gray}{\g/100}
        \fill[cc] (\px,\py) rectangle (\px+0.125,\py+0.125);
      }
    }
    \draw[thick] (0,0) rectangle (3,3);
    \node[below,font=\small] at (1.5,-0.1) {\lab};
  \end{scope}
}
\end{tikzpicture}
\caption{Three single-frequency 2-D patterns of identical fineness but different
\emph{orientation}. In 1-D these would be the same frequency; in 2-D they are
distinct, because spatial frequency is a vector $(u,v)$ carrying a direction. This
extra degree of freedom is the heart of what is new in two dimensions.}
\label{ch30:fig-orientation}
\end{figure}

\begin{intuitionbox}
Here is the cleanest way to feel the new dimension. A 2-D sinusoid
$\cos\!\big(2\pi(u x + v y)\big)$ is a corrugated sheet of light and dark stripes.
The vector $(u,v)$ points \emph{perpendicular to the stripes}, in the direction
the pattern changes fastest; its length $\sqrt{u^2+v^2}$ is how many cycles per
unit length you cross moving that way. Rotate $(u,v)$ and the stripes rotate;
lengthen it and the stripes get finer. Every image is a sum of these corrugations
at all orientations and finenesses -- that sum is the 2-D Fourier transform, and
it is where we go next.
\end{intuitionbox}

\begin{goalbox}
Where do images fit in a \emph{wrist}-worn gesture controller, whose primary
sensors are an IMU and (for finger-level intent) EMG or force myography? They earn
their place in two ways. The first is as a \emph{development and ground-truth
tool}: during prototyping you point a camera at the hand and use computer vision
(e.g.\ a hand-pose estimator) to label what the fingers were actually doing, then
train the wrist sensors to predict those labels. The camera never ships on the
wrist -- it is the teacher, not the product. The second is as a \emph{conceptual
bridge}: a
spectrogram (the chapter on the STFT) is literally a 2-D signal -- time across,
frequency up, magnitude as brightness -- so every tool here (histograms, contrast,
point operations, and soon 2-D filtering) applies directly to reading and cleaning
spectrograms of IMU and EMG data. The image chapters quietly upgrade your ability
to process the wrist signals themselves.
\end{goalbox}

\begin{focusbox}
Be clear-eyed about the limit. A camera \emph{can} see fingers and recover detailed
hand pose, but a camera is not wrist-wearable in any practical, all-day,
any-lighting sense. The wrist sensors that \emph{are} wearable -- IMU, EMG, FMG --
cannot themselves form an image; they give time series, not pixels. So vision is a
lab instrument here, not the deliverable. We will return to what the wearable
sensors truly can and cannot do in the hand-sensing chapter; the takeaway for now
is that ``just put a camera on it'' is not an answer to wrist-worn finger sensing.
\end{focusbox}

\section{A worked example}
\label{ch30:worked}

Let us run the numbers on a tiny image so the formulas stop being abstract. Take
the $3\times 3$ grayscale patch (8-bit)
\[
f = \begin{bmatrix} 30 & 60 & 60 \\ 60 & 90 & 90 \\ 90 & 120 & 150 \end{bmatrix}.
\]
There are $MN=9$ pixels. The histogram counts: one pixel at $30$, three at $60$,
three at $90$, one at $120$, one at $150$. Everything else is zero.

\textbf{Mean (brightness).}
$\bar f = \tfrac{1}{9}(30+60+60+60+90+90+90+120+150)=\tfrac{750}{9}\approx 83.3$.

\textbf{Min/max stretch.} Here $f_{\min}=30$, $f_{\max}=150$, so by
\eqref{ch30:stretch} each pixel maps to $g=255\,(f-30)/120$. For example
$60\mapsto 255\cdot 30/120 = 63.75\to 64$, and $90\mapsto 255\cdot 60/120=127.5\to
128$. The darkest pixel becomes $0$, the brightest $255$: full contrast.

\textbf{Equalization.} Build the normalized histogram and its CDF over the five
occupied levels:
\[
\begin{array}{c|ccccc}
k & 30 & 60 & 90 & 120 & 150\\ \hline
p[k] & 1/9 & 3/9 & 3/9 & 1/9 & 1/9\\
C[k] & 1/9 & 4/9 & 7/9 & 8/9 & 9/9
\end{array}
\]
Map each level by $k\mapsto \floor{255\,C[k]}$: the five levels become
$\floor{255/9}=28$, $\floor{255\cdot4/9}=113$, $\floor{255\cdot7/9}=198$,
$\floor{255\cdot8/9}=226$, and $255$. Notice the crowded middle levels ($60$ and
$90$, three pixels each) are pushed far apart -- from $30$ apart to $85$ apart --
buying contrast exactly where most pixels live, which is the whole point of
equalization.

\begin{walkbox}
Why did equalization spread $60$ and $90$ so much further than the stretch did?
Because the stretch is \emph{linear} -- it spaces output levels evenly with input
value -- while equalization spaces them by \emph{population}. The jump from $60$ to
$90$ crosses three pixels' worth of probability ($3/9$), so it earns three-ninths
of the output range. Sparse regions of the histogram get little output range;
dense regions get a lot. That is precisely ``allocate contrast where the pixels
are.''
\end{walkbox}

\begin{buildbox}
On an embedded target (Arduino-class MCU or a Raspberry Pi) point operations are
the cheapest image processing you can do, which matters when you are also running
sensor DSP. Three notes. (1)~Precompute the transfer curve into a
\texttt{uint8\_t lut[256]} once; the per-pixel cost is then a single array read --
no float math in the inner loop. (2)~Saturate, do not wrap: cast to a wider
integer, clamp to $[0,255]$, then store. (3)~If you only need brightness for
analysis (not display), skip the RGB-to-gray matrix multiply per pixel and use the
cheap integer approximation $Y\approx(77R+150G+29B)\!\gg\!8$, which is
\eqref{ch30:luma} scaled to $256$ and bit-shifted -- accurate to a level or two
and far faster on a chip with no FPU.
\end{buildbox}

\begin{lstlisting}[language=C,caption={A point operation as a precomputed LUT on an MCU. No floats in the inner loop.},captionpos=b]
#include <stdint.h>
#include <math.h>

static uint8_t lut[256];

void build_gamma_lut(float gamma) {        /* gamma < 1 brightens */
    for (int k = 0; k < 256; k++) {
        float v = 255.0f * powf(k / 255.0f, gamma);
        int   q = (int)(v + 0.5f);         /* round */
        lut[k] = (q < 0) ? 0 : (q > 255 ? 255 : (uint8_t)q);  /* saturate */
    }
}

void apply_lut(uint8_t *img, int n) {      /* n = width*height pixels */
    for (int i = 0; i < n; i++)
        img[i] = lut[img[i]];              /* one table read per pixel */
}
\end{lstlisting}

\section{Summary}
\label{ch30:summary}

An image is a signal that lives on a plane: a function $f(x,y)$ the world hands us
continuously, which a camera samples on a grid and quantizes into an array
$f[m,n]$. Each grid value is a pixel -- a \emph{sample}, not a square. A grayscale
pixel is one intensity; a color pixel is a 3-vector in a color space, and pulling
out the luma channel $Y$ \eqref{ch30:luma} separates brightness from color the way
the eye does. Two-dimensional sampling obeys the same Nyquist law as 1-D, applied
per axis, with aliasing showing up as moir\'e and jaggies. The histogram is the
value distribution -- a probability distribution over intensity -- and its first
two moments \emph{are} brightness and contrast. Point operations remap each pixel
through a single transfer curve $T$, implementable as a fast lookup table:
brightness shifts, contrast stretches, gamma, negatives, and thresholds are all
just curve shapes. The genuinely new ingredient in two dimensions is
\emph{direction}: spatial frequency is a vector with an orientation, which is what
makes the next chapter's 2-D Fourier transform a picture rather than a plot.

\section{Exercises}
\label{ch30:exercises}

\begin{exercise}
\label{ch30:ex-vector}
An RGB image is $1280\times 720$ with 8 bits per channel. (a)~How many pixels does
it have? (b)~How many total bytes, uncompressed? (c)~Viewed as a single vector,
what is its dimension? (d)~If you convert to single-channel luma, how many bytes
and what dimension?
\end{exercise}

\begin{exercise}
\label{ch30:ex-luma}
A pixel has $(R,G,B)=(200,100,50)$. (a)~Compute its luma $Y$ from
\eqref{ch30:luma}. (b)~Compute the fast integer approximation
$Y\approx(77R+150G+29B)\gg 8$ and compare. (c)~Explain in one sentence why green
dominates the weights.
\end{exercise}

\begin{exercise}
\label{ch30:ex-affine}
You want to map 8-bit intensities so that input $40$ goes to output $0$ and input
$210$ goes to output $255$, linearly. (a)~Find the gain $a$ and offset $b$ in
$g=af+b$. (b)~What output does input $125$ produce? (c)~What happens to inputs
below $40$ or above $210$, and what must you do about it?
\end{exercise}

\begin{exercise}
\label{ch30:ex-hist}
Consider the $2\times 2$ image $\big[\begin{smallmatrix}10&10\\200&200\end{smallmatrix}\big]$.
(a)~Write its histogram. (b)~Compute the mean and variance of intensity. (c)~Is
this high or low contrast, and how does the variance tell you?
\end{exercise}

\begin{exercise}
\label{ch30:ex-equalize}
For the same image as Exercise~\ref{ch30:ex-hist}, perform histogram equalization
using $k\mapsto\floor{255\,C[k]}$. Report the new pixel values. Did contrast
increase? Why is equalization nearly useless on a two-level image, and what kind
of histogram makes it most effective?
\end{exercise}

\begin{exercise}
\label{ch30:ex-gamma}
(a)~Sketch the transfer curves for $\gamma=0.5$ and $\gamma=2.0$ on the same axes
as the identity. (b)~For an 8-bit input of $64$, compute the output under each
gamma using \eqref{ch30:gamma-eq}. (c)~Which gamma brightens the shadows, and
explain how the slope of the curve near $f=0$ tells you so.
\end{exercise}

\begin{exercise}
\label{ch30:ex-gammablend}
Two pixels are gamma-encoded ($\gamma_{\text{enc}}=1/2.2$) with values $0$ and
$255$. (a)~Naively average the encoded values. (b)~Instead, decode each to linear
light, average, then re-encode. (c)~Compare the two results and explain why the
naive average looks too dark.
\end{exercise}

\begin{exercise}
\label{ch30:ex-threshold}
You photograph black text on white paper. (a)~Sketch the expected histogram.
(b)~Where would you place a threshold $\theta$ and why? (c)~Describe in words how
Otsu's method would choose $\theta$ from the histogram automatically, and connect
it to the idea of separating a distribution into two clusters.
\end{exercise}

\begin{exercise}
\label{ch30:ex-alias}
A camera samples at one pixel per millimetre and photographs a fabric with
$0.7$\,mm-wide alternating stripes (period $1.4$\,mm). (a)~What spatial frequency
is the fabric, in cycles per mm? (b)~What is the Nyquist limit for this sampling?
(c)~Will you see aliasing, and what does the moir\'e pattern's apparent frequency
become? (d)~What hardware fix do camera makers use?
\end{exercise}

\begin{exercise}
\label{ch30:ex-orientation}
Write the 2-D sinusoid $\cos\!\big(2\pi(u x+v y)\big)$ for $(u,v)=(0,4)$,
$(4,0)$, and $(3,3)$ cycles per unit. For each, state the stripe orientation, the
direction of fastest change, and the fineness $\sqrt{u^2+v^2}$. Which two have the
same fineness, and why are all three nonetheless distinct signals?
\end{exercise}

\begin{exercise}
\label{ch30:ex-spectrogram}
A spectrogram of forearm-IMU gyro data is being treated as an image (time on the
horizontal axis, frequency on the vertical, magnitude as brightness). (a)~What
would a histogram of this image tell you about the recording? (b)~Which point
operation would you apply to make faint high-frequency activity visible, and what
shape transfer curve does that need? (c)~Explain why this is the \emph{same}
operation a radiologist's window/level control performs.
\end{exercise}

\begin{exercise}
\label{ch30:ex-camera}
A teammate proposes mounting a small camera on the wristband so the gesture
controller can see the fingers directly. Give two concrete reasons this is a poor
choice for a wearable product, and explain the role a camera \emph{should} play in
developing the system (hint: ground truth). Name the sensor modalities that are
actually wrist-wearable for finger-level intent.
\end{exercise}
